\addcontentsline{toc}{chapter}{Abstract}\vspace{-1cm}
\begin{tikzpicture}[remember picture, overlay]
  \draw[line width = 4pt] ($(current page.north west) + (0.75in,-0.75in)$) rectangle ($(current page.south east) + (-0.75in,0.75in)$);
\end{tikzpicture}

Intrusion Detection Systems are an essential component of modern network security, but conventional software-based machine learning approaches depend on repeated numerical computation and continuous monitoring. Such approaches can provide high detection accuracy, yet their computational style is not naturally suited for ultra-low-power, always-on edge security devices. Neuromorphic computing provides an alternative direction by representing information as sparse spike events and processing those events through neuron-inspired circuits.

This interdisciplinary project presents the design foundation and implementation workflow for a Neuromorphic Intrusion Detection System based on Spiking Neural Networks and Leaky Integrate-and-Fire neuron behavior. The project combines Information Science and Engineering aspects such as data preprocessing, dimensionality reduction, spike encoding, and classification with Electronics and Communication Engineering aspects such as analog neuron modeling and Cadence-compatible waveform generation. The UNSW-NB15 network intrusion dataset is processed, normalized, reduced using Principal Component Analysis, and mapped into spike-based voltage waveforms suitable for circuit-level evaluation.

The implemented software pipeline generates analog-compatible Piecewise Linear voltage source files from processed network traffic features. These files serve as direct stimulus inputs for Leaky Integrate-and-Fire neuron circuits designed in SPICE/Cadence environments. The report documents the motivation, theoretical fundamentals, system architecture, design methodology, and implementation details completed for the project. The present report is limited to Chapters 1 through 4, focusing on the project foundation, design, and implementation workflow. Detailed result analysis and concluding validation are intentionally excluded from this document scope.

\pagebreak
